% Modernised reading — fonds Alexandre Grothendieck, shelfmark 151, batch 1
% (pages 1-20).
%
% This is an INTERPRETATION, not a transcription. It re-reads the pages in
% today's notation and vocabulary, states the hypotheses the manuscript leaves
% implicit, and drops the critical apparatus so that the mathematics reads
% straight through. Everything the brackets and underlines carried — a doubtful
% reading, a notation he used differently, a missing passage — moves into a
% footnote.
%
% It is held to being mathematically correct as it stands, not merely faithful
% to the page. Where the manuscript is loose or mistaken, the true statement is
% what appears here, and the footnote says what the page has. In any dispute of
% substance the transcription governs: whoever cites Grothendieck must cite
% batch-01.fr.tex.
%
% Scope: the folder was read WHOLE before any of this was written — all four
% batches, pages 1-74, in one pass — and one file is written per batch, this
% being the first. That matters here more than anywhere else in the folder:
% pages 2-8 and 16-20 are working sheets whose conventions are only fixed on
% page 21, in the next batch, and the sign conventions used below are read back
% from there.
%
% Pass: Opus 5 (claude-opus-5), 2026-08-16 — first pass, unchecked against the
% pages by a human.
\documentclass[11pt,a4paper]{article}
% Le préambule commun à toutes les transcriptions du fonds Grothendieck.
%
% Il tient en une idée : **l'incertitude se compose, elle ne se résout pas.**
% Une transcription qui lisse un mot douteux a détruit la seule chose pour
% laquelle on la faisait — un lecteur doit pouvoir savoir, à chaque endroit,
% ce qui était sur le feuillet et ce qui a été ajouté. Les six macros qui
% suivent sont donc l'appareil critique, et non de la décoration.
%
% Les mêmes commandes sont reconnues par `scripts/render.mjs`, qui produit la
% vue de lecture du site. Une macro ajoutée ici doit l'être aussi là-bas,
% faute de quoi le rendu échouera bruyamment — ce qui est voulu : un rendu
% silencieusement faux est pire qu'un rendu absent.

\ProvidesPackage{grothendieck}[2026/08/08 Transcriptions du fonds Grothendieck]

\usepackage{amsmath,amssymb,amsthm}
\usepackage{mathrsfs}
\usepackage{stmaryrd}
% Les diagrammes commutatifs sont partout dans ce fonds : c'est la langue
% même de la géométrie algébrique de Grothendieck, et une transcription qui
% les rendrait en prose ne transcrirait rien.
\usepackage{tikz-cd}
% Français par défaut — c'est la langue des feuillets — mais l'anglais est
% chargé aussi : la traduction et le résumé sont des documents anglais, et la
% typographie française (espaces avant les deux-points) y serait une faute.
% Ces documents basculent par \selectlanguage{english} en tête de corps.
\usepackage[english,french]{babel}
\usepackage{fontspec}
\usepackage{geometry}
\geometry{a4paper,margin=25mm,marginparwidth=28mm}
\usepackage{marginnote}
\usepackage{xcolor}
% ulem, not soul. `soul`'s \st and \ul are built on a character-by-character
% scan that XeLaTeX's font machinery breaks: the document fails with an
% undefined control sequence rather than degrading. ulem does the same two jobs
% and is Unicode-engine safe. `normalem` stops it hijacking \emph, which the
% transcriptions use for Grothendieck's own emphasis.
\usepackage[normalem]{ulem}
\usepackage{hyperref}
\hypersetup{colorlinks=true,linkcolor=black,urlcolor=black,pdfborder={0 0 0}}

% --- Métadonnées du lot -----------------------------------------------------
% Reprises telles quelles par le script de rendu, qui les affiche en tête de
% la vue de lecture. Elles fixent ce que le fichier prétend couvrir : sans
% elles, une transcription flotte sans point d'attache dans un fonds de seize
% mille feuillets.

\newcommand{\folder}[1]{\gdef\grothFolder{#1}}
\newcommand{\batch}[1]{\gdef\grothBatch{#1}}
\newcommand{\leaves}[2]{\gdef\grothFirst{#1}\gdef\grothLast{#2}}
\newcommand{\foldertitle}[1]{\gdef\grothFoldertitle{#1}}
\gdef\grothFolder{?}\gdef\grothBatch{?}\gdef\grothFirst{?}\gdef\grothLast{?}\gdef\grothFoldertitle{}

% --- Le résumé --------------------------------------------------------------
%
% Il ouvre la lecture modernisée, sous le titre « Résumé », et il est composé
% en petit. Ce n'est pas un ornement : c'est le seul passage du document qui
% ne soit pas des mathématiques, et un lecteur doit pouvoir le reconnaître
% avant de l'avoir lu — puis le sauter s'il connaît déjà le sujet, ou s'y
% arrêter s'il ne le connaît pas. Le filet qui le suit marque la reprise du
% corps.

\newenvironment{resume}
  {\small\setlength{\parskip}{0.45em}}
  {\par\medskip\hrule\medskip}

% --- Mots-clés ---------------------------------------------------------------
%
% En anglais, en clôture du résumé de la lecture modernisée : le vocabulaire
% moderne sous lequel chercher ce que les feuillets construisent. Le manifeste
% du site (`npm run manifest`) les extrait pour étiqueter la cote entière —
% cette ligne est la source unique des tags, il n'y a pas de fichier à côté.
\newcommand{\keywords}[1]{\par\smallskip\noindent
  {\footnotesize\textsc{Keywords} --- \textit{#1}}\par}

% --- L'appareil critique ----------------------------------------------------

% Le numéro de feuillet, en marge. C'est l'ancre qui permet de régler un
% désaccord sur une formule en regardant *un* feuillet plutôt qu'une chemise
% de six cents. Le site s'en sert aussi pour tourner les pages du fac-similé
% quand on fait défiler la transcription.
\newcommand{\leaf}[1]{%
  \marginnote{\footnotesize\color{gray!70}\textsf{#1}}%
  \label{leaf:#1}\ignorespaces}

% Illisible : on ne devine pas. Un mot inventé qui se lit comme les autres est
% le pire résultat possible.
\newcommand{\ill}{\textcolor{red!60!black}{[\,\ldots\,]}}

% Lecture douteuse : le mot est donné, et signalé comme incertain.
\newcommand{\uncertain}[1]{\uline{#1}}

% Ajout de l'éditeur — une lettre manquante, un mot sous-entendu.
\newcommand{\add}[1]{\textcolor{blue!50!black}{[#1]}}

% Biffé par Grothendieck lui-même. Ce qu'il a rayé fait partie du document :
% une rature dit souvent où la pensée a changé de direction.
\newcommand{\struck}[1]{\sout{#1}}

% Note du transcripteur, sur l'état matériel du feuillet.
\newcommand{\note}[1]{\par\smallskip{\small\color{gray!60!black}\textit{[#1]}}\par\smallskip}

% Note marginale de Grothendieck, distincte du corps du texte.
\newcommand{\marginal}[1]{\par\smallskip\noindent\hspace{1em}%
  {\small\color{gray!60!black}#1}\par\smallskip}

% --- Titre ------------------------------------------------------------------

\AtBeginDocument{%
  \begin{center}
    {\large\bfseries Fonds Alexandre Grothendieck --- cote n\textsuperscript{o}~\grothFolder}\\[2pt]
    {\itshape\grothFoldertitle}\\[4pt]
    {\small feuillets \grothFirst--\grothLast\ (lot \grothBatch)}\\[2pt]
    {\footnotesize\color{gray!60!black}Transcription établie d'après le fac-similé numérisé
     par l'Université de Montpellier.\\
     Les lectures incertaines sont soulignées, les illisibles notées [\,\ldots\,],
     les ajouts entre crochets.}
  \end{center}
  \vspace{1em}}

\endinput


\folder{151}
\batch{1}
\pages{1}{20}
\dating{[à partir de 1981-à partir de 1990]}
\watermark{Édition de démonstration\\interprétation personnelle de l'œuvre}
\foldertitle{Topos stratifié (1981) : notes manuscrites (s.d.) — lecture modernisée}

\begin{document}

\section*{Résumé}

\begin{resume}
Un espace stratifié est un espace découpé en morceaux : une surface et les
courbes qu'elle porte, une variété algébrique et son lieu singulier, un espace
de modules et son compactifié. Chaque morceau, pris seul, est lisse ; les
singularités n'apparaissent qu'aux jointures. Ce dossier poursuit une question
unique : \emph{si l'on connaît les morceaux, et la manière dont chacun
s'approche des autres, connaît-on l'espace ?}

La réponse naïve — recoller les morceaux le long de leurs intersections — ne
peut pas suffire, et pour une raison qu'on voit sur un dessin. Prenez un cône,
et pour morceaux la pointe et le reste. Leur intersection est vide : la pointe
est fermée, le reste est ouvert, et rien dans ces deux données ne dit que le
cône se referme sur sa pointe. Il faut donc porter, en plus des morceaux, un
enregistrement de la façon dont l'un s'approche de l'autre — un voisinage de la
pointe, et ce qu'il en reste quand on ôte la pointe elle-même. C'est ce que
Grothendieck appelle un \emph{tube}, et tout le vocabulaire du dossier tourne
autour de ce mot : le tube $\mathcal{V}_{i,j}$ d'une strate dans une autre, le
tube épointé $\mathcal{V}^{*}_{i,j}$, et les tubes mixtes qui les interpolent.

L'idée qui arrive ensuite est celle d'un diagramme d'indices. Les morceaux sont
indexés par un ensemble ordonné $I$ — la relation $i \leq j$ voulant dire que la
$i$-ième strate est dans l'adhérence de la $j$-ième — et à cet ensemble ordonné
on associe un diagramme d'espaces : les strates, les tubes, les tubes épointés,
et entre eux trois espèces de flèches. Le théorème que vise le dossier dit que
l'espace tout entier se retrouve comme recollement de ce diagramme. Il est
énoncé à la page 68, sous le nom de \emph{théorème de recollement des tubes}, et
sa démonstration n'est pas dans ces pages.

Sur le chemin, un piège, et il est le moment le plus instructif du dossier. On
peut être tenté de ne retenir des morceaux que leur groupe fondamental, et des
jointures que les homomorphismes correspondants. Grothendieck le fait, à la
section 1, puis calcule ce que ce niveau de données donne pour la sphère
$\mathbb{P}^{1}_{\mathbb{C}} \simeq S^{2}$ vue comme le plan complexe plus un
point à l'infini — et trouve un objet homotopiquement trivial, alors que $S^{2}$
ne l'est pas. Les morceaux sont pourtant les plus simples possibles. La
conclusion qu'il en tire, aux pages 40 à 42, est qu'il faut monter d'un cran :
retenir non seulement les groupes mais les gerbes, et il retrouve alors la bonne
valeur $H^{2}(S^{2},A) \simeq A$ — « un magnifique succès ! », écrit-il.

Où cela mène-t-il ? À ce qu'on appelle aujourd'hui la théorie de l'homotopie
stratifiée : catégorie des chemins sortants, faisceaux constructibles sur un
espace stratifié, colimite homotopique sur le poset des strates. Le programme
que ces pages formulent — exprimer le type d'homotopie d'un espace stratifié
dont les pièces sont des $K(\pi,1)$ à partir des groupoïdes fondamentaux de ces
pièces — est celui-là même.

Les vingt pages du présent lot ne sont pas ce texte. Elles en portent le plan,
à la page 1, et trois séries de notes de travail qui l'entourent sans le citer.
La première, aux pages 2 à 8, est un calcul de relations dans un système
$E_{0}, E_{1}, E_{2}$ à trois étages, muni des permutations des coordonnées :
c'est la forme abstraite du \emph{nerf de Čech} d'un recouvrement, écrite au
niveau où l'on peut la vérifier à la main. La deuxième, aux pages 10 à 14, est
un texte suivi, « Hyperrecouvrements », qui demande ce que doit être un
hyperrecouvrement d'un topos indexé par une petite catégorie quelconque. La
troisième, aux pages 16 à 20, reprend le premier calcul en toutes dimensions et
le dualise.

Ces trois pièces posent la même question sous trois habits : \emph{quand un
système de pièces et de recouvrements multiples détermine-t-il le tout ?} La
condition y prend chaque fois la même forme — les recouvrements multiples se
surjectent sur les produits — et c'est cette forme qui reviendra, transposée aux
strates, dans le corps de « Structures stratifiées ». Le lien est thématique et
non textuel : aucune de ces pages ne renvoie aux autres.

\keywords{hypercovering, Čech nerve, flat functor, Diaconescu theorem, symmetric simplicial object}
\end{resume}

\pagerange{1}{1}
\section*{Le plan de « Structures stratifiées »}

La première page est le sommaire d'un texte en quatre sections, avec la
pagination de l'auteur en marge :

\begin{enumerate}
  \item La situation type — préliminaires d'un programme (p.~1 de l'auteur).
  \item Stratifications globales : le formalisme $\tfrac{1}{2}$ simplicial
    (sans tubes) (p.~10).
  \item Introduction des tubes (p.~15).
  \item Topos canoniques associés à une stratification globale (p.~23).
\end{enumerate}

Ce sommaire est tenu. Le texte lui-même commence à la page 26 du fonds, dans le
lot suivant, et court sans interruption jusqu'à la dernière page du dossier : la
section~1 occupe les pages 26 à 43, la section~2 les pages 44 à 55, la
section~3 les pages 56 à 68, la section~4 les pages 69 à 74, où le dossier
s'arrête au milieu d'une construction.\footnote{La pagination de l'auteur et
celle des archivistes ne coïncident pas : l'auteur numérote ses feuillets, dont
chacun porte deux pages du fonds. Les renvois du sommaire sont à sa
numérotation ; ceux du présent commentaire, comme partout dans cette édition,
sont à celle des archivistes.}

\pagerange{2}{8}
\section*{Le nerf de Čech tronqué : le système $E_{0}, E_{1}, E_{2}$}

Quatre feuilles de travail, numérotées 1 à 4 par l'auteur, calculent dans un
système à trois étages. Le modèle qui les rend toutes transparentes n'est pas
écrit sur la page mais s'y lit dans les formules en coordonnées : partant d'un
ensemble $F$, on pose
\[
  E_{0} = F, \qquad E_{1} = F \times F, \qquad E_{2} = F \times F \times F,
\]
les $p$ oubliant une coordonnée, les $\delta$ en répétant une, et
$\mathfrak{S}_{n+1}$ permutant les facteurs de $E_{n}$ :
\[
  p_{0}(x,y) = x, \quad p_{1}(x,y) = y, \quad \delta_{0}(x) = (x,x),
  \quad \sigma(x,y) = (y,x),
\]
\[
  p_{12}(x,y,z) = (x,y), \quad p_{23}(x,y,z) = (y,z), \quad
  p_{31}(x,y,z) = (z,x),
\]
\[
  \delta_{1}(x,y) = (y,x,x), \quad \delta_{2}(x,y) = (x,y,x), \quad
  \delta_{3}(x,y) = (x,x,y), \quad \rho(x,y,z) = (y,z,x).
\]
C'est le \emph{nerf de Čech} de l'application $F \to \{*\}$, tronqué au
deuxième étage. Ce que les quatre feuilles cherchent est sa présentation : quel
est le plus petit système de données et de relations dont ce modèle soit
l'exemple libre.

\subsection*{La convention sur les transpositions}

Une convention doit être fixée avant tout, car deux relations de ces pages sont
fausses sans elle et vraies avec. Les transpositions $\sigma_{1}, \sigma_{2},
\sigma_{3}$ de $\mathfrak{S}_{3}$ sont indexées par la lettre qu'elles
\emph{fixent} :
\[
  \sigma_{1}(x,y,z) = (x,z,y), \qquad
  \sigma_{2}(x,y,z) = (z,y,x), \qquad
  \sigma_{3}(x,y,z) = (y,x,z).
\]
C'est la convention qu'exigent les relations d'équivariance
\[
  \sigma\, p_{12} = p_{12}\,\sigma_{3}, \qquad
  \sigma\, p_{23} = p_{23}\,\sigma_{1}, \qquad
  \sigma\, p_{31} = p_{31}\,\sigma_{2},
\]
qui sont vraies dans le modèle et sont écrites telles quelles sur la
page.\footnote{La convention n'est explicitée nulle part dans ce lot ; elle
l'est à la page 21, dans le lot suivant, où les trois transpositions sont
données en coordonnées. Nous la lisons donc à rebours, ce que seule une lecture
du dossier entier permet.}

\subsection*{Le système et ses relations}

La donnée minimale est celle de trois ensembles $E_{0}$, $(E_{1},
\mathfrak{S}_{2})$, $(E_{2}, \mathfrak{S}_{3})$ et de quatre applications
$p_{0} \colon E_{1} \to E_{0}$, $\delta_{0} \colon E_{0} \to E_{1}$,
$p_{12} \colon E_{2} \to E_{1}$, $\delta_{1} \colon E_{1} \to E_{2}$ ; tout le
reste s'en déduit par les groupes :
\[
  p_{1} = p_{0}\sigma, \qquad
  p_{23} = p_{12}\rho, \qquad p_{31} = p_{12}\rho^{2}, \qquad
  \delta_{3} = \rho\,\delta_{1}, \qquad \delta_{2} = \rho^{2}\delta_{1}.
\]
Les relations que la page encadre, et qui sont toutes vraies dans le modèle,
sont
\[
  p_{0}\delta_{0} = \mathrm{id}_{E_{0}}, \qquad \sigma\delta_{0} = \delta_{0},
\]
\[
  p_{12}\delta_{1} = \sigma, \qquad
  p_{12}\delta_{3} = \delta_{0}p_{0}, \qquad
  p_{12}\delta_{2} = \mathrm{id}_{E_{1}},
\]
\[
  \sigma_{1}\delta_{1} = \delta_{1}, \qquad
  p_{12}\sigma_{3} = \sigma\,p_{12}.
\]
Les trois du milieu sont les trois valeurs distinctes que prend $p_{12}$ sur
les trois faces $\delta_{1}, \delta_{2}, \delta_{3}$ : une involution, un
idempotent, l'identité. Elles sont, à la lettre, trois des identités
simpliciales.\footnote{La page 2 écrit $\sigma_{3}\delta_{1} = \delta_{1}$ ; la
page 4 écrit $\sigma_{1}\delta_{1} = \delta_{1}$. Sous la convention ci-dessus,
c'est la seconde qui est vraie : $\delta_{1}(x,y) = (y,x,x)$ est fixé par la
transposition des deux dernières lettres, c'est-à-dire par $\sigma_{1}$. Le
stabilisateur de $\delta_{1}(\xi)$ dans $\mathfrak{S}_{3}$ est exactement
$\{1,\sigma_{1}\}$ dès que $\xi$ n'est pas diagonal, et c'est ce fait qui sert
au lot suivant.}

Deux conséquences méritent d'être isolées, parce qu'elles servent plus loin.
La première est que $\delta_{1}\delta_{0}$ est invariante sous $\mathfrak{S}_{3}$
tout entier :
\[
  g\,(\delta_{1}\delta_{0}) = \delta_{1}\delta_{0}
  \qquad (g \in \mathfrak{S}_{3}),
\]
ce qui est immédiat sur le modèle, où $\delta_{1}\delta_{0}(x) = (x,x,x)$. La
seconde est que $p_{0}p_{12}$ est invariante à droite sous $\sigma_{1}$ :
\[
  p_{0}\,p_{12}\,\sigma_{1} = p_{0}\,p_{12} = p_{0}\,p_{13},
\]
$\sigma_{1}$ ne touchant pas la première lettre.\footnote{La page 4 écrit cette
dernière relation avec $\sigma_{3}$ et nomme le membre de droite $p_{1}$ ; les
deux sont fautifs, et de la même façon que la relation
$\sigma_{3}\delta_{1} = \delta_{1}$ de la page 2. Sous la convention fixée,
$p_{0}p_{12}\sigma_{3}$ vaut $p_{1}p_{12}$ et non $p_{0}p_{12}$. La page 6
écrit la relation correcte, avec $\sigma_{1}$.}

\subsection*{Une inégalité qui n'en est pas une}

La page 2 porte, barrée et flanquée de deux points d'interrogation,
\[
  (p_{12}\rho)\,\delta_{1} \neq p_{23}\,(\rho\delta_{1}) \qquad ??
\]
L'inégalité est vraie et n'a rien d'inquiétant. Comme $p_{12}\rho = p_{23}$, le
membre de gauche est $p_{23}\delta_{1} = \delta_{0}p_{0}$ ; le membre de droite
est $p_{23}\delta_{3} = \mathrm{id}_{E_{1}}$. Ce ne sont pas deux
parenthésages du même composé — le second n'est pas
$p_{12}(\rho\delta_{1})$ — et la composition n'est pas en cause. Les deux
valeurs obtenues sont précisément les deux premières lignes de la boîte que la
page écrit juste en dessous.\footnote{Le doute vient des accolades dont l'auteur
surcharge la ligne, qui renomment $\rho\delta_{1}$ en $\delta_{3}$ et
$\rho^{2}\delta_{1}$ en $\delta_{2}$ : un même symbole y sert de nom et de
composé.}

\subsection*{Le cas pointé, et la lecture cocyclique}

Si $E_{0}$ est réduit à un point, $E_{1}$ et $E_{2}$ deviennent des ensembles
pointés par $e_{1} = \delta_{0}(*)$ et $e_{2} = \delta_{1}(e_{1})$, les groupes
respectent les pointages, et $\delta_{0}p_{0}$ devient l'application constante.
Les relations se réduisent à cinq :
\[
  p_{12}\delta_{1} = \sigma, \qquad
  p_{12}\delta_{3} = e, \qquad
  p_{12}\delta_{2} = \mathrm{id}_{E_{1}}, \qquad
  \sigma_{1}\delta_{1} = \delta_{1}, \qquad
  p_{12}\sigma_{3} = \sigma\,p_{12},
\]
où $e$ est l'application constante. La page 8 les récrit multiplicativement,
et c'est le seul endroit du lot où le système soit lu comme une présentation :
\[
  \delta_{0}p_{0}(x) = 1, \qquad
  \sigma(x)\cdot x = 1, \qquad
  p_{12}(z)\,p_{23}(z)\,p_{31}(z) = 1 \quad (z \in E_{2}).
\]
Ce sont, dans l'ordre, la normalisation, l'inversion et la condition de
cocycle. Autrement dit : $E_{1}$ joue le rôle des données de transition d'un
recouvrement, $\sigma$ celui du passage à l'inverse, $\delta_{0}$ celui des
identités, et la troisième relation est la condition de cocycle de Čech sur les
intersections triples.\footnote{La lecture est celle que la page impose par sa
notation multiplicative ; elle n'y est pas nommée. On ajoutera qu'un tel système
avec ses groupes symétriques est, en langue d'aujourd'hui, un \emph{objet
simplicial symétrique} tronqué au deuxième étage — un foncteur sur la catégorie
des ensembles finis non vides à au plus trois éléments — et que le modèle
$E_{n} = F^{n+1}$ en est l'exemple universel.}

Une remarque de marge de la page 8 vaut d'être conservée, car c'est un critère
et non une étape : $\delta_{1}(x)$ est fixé par $\rho$ — donc par
$\mathfrak{S}_{3}$ tout entier — si et seulement si $x = \delta_{0}p_{0}(x)$,
c'est-à-dire si et seulement si $x$ est dégénéré.

\pagerange{10}{14}
\section*{Hyperrecouvrements}

Cinq pages d'un texte suivi, paginées 1 à 5 par l'auteur, demandent ce que doit
être un hyperrecouvrement d'un topos $X$ \emph{indexé par une petite catégorie
$A$ quelconque}, là où la notion reçue est indexée par un objet
simplicial.\footnote{Les hyperrecouvrements de SGA 4, exposé V, et ceux
d'Artin--Mazur sont indexés par des objets simpliciaux ; le texte cite d'ailleurs
Artin--Mazur quelques pages plus loin, à un autre propos. Le passage à une
catégorie d'indices arbitraire est ce que ces pages proposent, et la
reformulation qu'elles atteignent est celle qui rend le passage indolore.}
La donnée est un foncteur
\[
  \varphi \colon A \longrightarrow X, \qquad \alpha \longmapsto U_{\alpha},
\]
et la question est : quelles conditions sur $\varphi$ ?

\subsection*{Les quatre axiomes}

Ils sont posés dans cet ordre.

\begin{itemize}
\item \textbf{Hyp 1.} Les $U_{\alpha}$, $\alpha \in \mathrm{Ob}\,A$, recouvrent
l'objet final $e_{X}$.
\item \textbf{Hyp 1${}'$.} Pour tout $\alpha_{0}$, le foncteur induit
$A/\alpha_{0} \to X/U_{\alpha_{0}}$ satisfait Hyp 1 : les
$U_{\alpha} \to U_{\alpha_{0}}$ indexées par les flèches $\alpha \to
\alpha_{0}$ forment une famille couvrante.
\item \textbf{Hyp 2.} Pour $\alpha, \beta$, la famille des
$U_{\gamma} \to U_{\alpha} \times U_{\beta}$, indexée par les objets de
$A/\alpha\times\beta$ — le produit étant pris dans $\widehat{A}$ —, est
couvrante.
\item \textbf{Hyp 2${}'$.} La version relative : pour tout diagramme
$\alpha_{1} \to \alpha_{0} \leftarrow \alpha_{2}$, la famille des
$U_{\gamma} \to U_{\alpha_{1}} \times_{U_{\alpha_{0}}} U_{\alpha_{2}}$,
indexée par les objets de $A/\alpha_{1}\times_{\alpha_{0}}\alpha_{2}$, est
couvrante.
\end{itemize}

La version à $n$ facteurs de Hyp 2 ne coûte rien : elle se déduit du cas
$n = 1$ par récurrence, puisque la composée de deux familles couvrantes est
couvrante, en écrivant la famille des
$U_{\gamma'} \to U_{\gamma} \times U_{\alpha_{n+1}} \to
(U_{\alpha_{0}} \times \cdots \times U_{\alpha_{n}}) \times
U_{\alpha_{n+1}}$.\footnote{Les deux étiquettes de ce diagramme sont
interverties sur la page relativement à la phrase qui le suit ; la
transcription le signale et ne corrige pas. Rien de mathématique n'en dépend :
la première flèche est celle du cran $n+1$, la seconde celle de l'hypothèse de
récurrence.}

\subsection*{La reformulation, et ce qu'elle dit}

Soit $\varphi_{!} \colon \widehat{A} \to X$ l'unique prolongement de $\varphi$
commutant aux petites colimites — l'extension de Kan à gauche le long de
Yoneda —, et soit
\[
  \varphi^{*} \colon X \longrightarrow \widehat{A}, \qquad
  \varphi^{*}(Y) = \bigl(\alpha \mapsto \mathrm{Hom}_{X}(\varphi\alpha, Y)\bigr)
\]
son adjoint à droite. Le premier est cocontinu, le second continu, et
$\varphi_{!} \dashv \varphi^{*}$.\footnote{L'étoile de $\varphi^{*}$ est ici
celle d'un adjoint à \emph{droite}, à rebours de l'usage : si $\varphi_{!}$
provient d'un morphisme de topos $f$, on a $\varphi_{!} = f^{*}$ et
$\varphi^{*} = f_{*}$. Le manuscrit écrit les deux notations à quelques lignes
l'une de l'autre, page 13 et page 17.}

Alors Hyp 1, Hyp 2, Hyp 2${}'$ sont exactement :
\begin{itemize}
\item a) $\varphi_{!}(e_{\widehat{A}}) \to e_{X}$ est épi ;
\item b) $\varphi_{!}(\alpha\times\beta) \to \varphi_{!}(\alpha) \times
\varphi_{!}(\beta)$ est épi, pour $\alpha,\beta \in \mathrm{Ob}\,A$ ;
\item c) $\varphi_{!}(\alpha\times_{\alpha_{0}}\beta) \to \varphi_{!}(\alpha)
\times_{\varphi_{!}(\alpha_{0})} \varphi_{!}(\beta)$ est épi, pour tout
diagramme $\alpha \to \alpha_{0} \leftarrow \beta$ dans $A$.
\end{itemize}
Et comme tout objet de $\widehat{A}$ est colimite de représentables,
$F = \varinjlim_{A/F}\alpha$, les conditions b) et c) valent alors pour des
$F, G$ quelconques de $\widehat{A}$, ce qui permet de tout ramasser en un
axiome unique :
\[
  \textbf{(Hyp)} \qquad
  \varphi_{!}\Bigl(\varprojlim_{I} F_{i}\Bigr) \longrightarrow
  \varprojlim_{I} \varphi_{!}(F_{i}) \quad \text{est épi},
\]
pour tout diagramme fini $(F_{i})_{i \in I}$ dans $\widehat{A}$ — et il suffit
de le postuler pour $I$ vide, discret à deux objets, et
$({\cdot} \to {\cdot} \leftarrow {\cdot})$, ce qui redonne respectivement
Hyp 1, Hyp 2, Hyp 2${}'$.\footnote{La démonstration que b) et c) sur les
représentables impliquent leur version générale est menée à la page 17, sur des
$\varinjlim$ indexées par $A/F$ et $A/G$ ; elle utilise que ces catégories
d'indices sont filtrantes deux à deux au sens voulu, ce que la page ne dit pas
mais que la construction du produit dans $\widehat{A}$ fournit.}

Ce que dit (Hyp) est net : \emph{$\varphi_{!}$ commute aux limites finies, à
un recouvrement près}. Si l'on y remplace « épi » par « iso », on obtient
l'exactitude à gauche de $\varphi_{!}$ ; le manuscrit relève lui-même ce cas
comme « particulièrement agréable » et note qu'il équivaut à la donnée d'un
morphisme de topos $f \colon X \to \widehat{A}$ avec $\varphi_{!} = f^{*}$.
C'est le théorème de Diaconescu, et ce que le manuscrit appelle le cas agréable
s'appelle aujourd'hui la \emph{platitude} du foncteur
$\varphi$.\footnote{R.~Diaconescu, \emph{Change of base for toposes with
generators}, 1975 : les morphismes de topos $X \to \widehat{A}$ correspondent
aux foncteurs plats $A \to X$, c'est-à-dire à ceux dont l'extension de Kan à
gauche est exacte à gauche. Le résultat est antérieur à ces pages ; rien
n'indique qu'elles s'y réfèrent, et le manuscrit énonce la correspondance sans
la nommer.}

Le texte se termine sur une question, laissée ouverte :

\begin{quote}
Mais cette hypothèse peut-elle être vérifiée, sans que $\varphi_{!}$ soit
exact ? Quand ?
\end{quote}

C'est exactement la question de l'écart entre la platitude et sa forme locale.
L'écart est réel — les deux conditions ne coïncident pas en général, et c'est
pour cela que la seconde a reçu un nom propre, \emph{covering flat} — mais le
dossier n'y revient pas.\footnote{La condition (Hyp) est, mot pour mot, celle
qu'on appelle aujourd'hui platitude relative à la topologie ; elle est ce qui
remplace la platitude dans l'énoncé de Diaconescu lorsque la catégorie
d'indices n'a pas les limites finies. Ces pages en donnent la formulation
correcte sans en tirer d'exemple séparant.}

\pagerange{16}{20}
\section*{Le système $(E_{n}, \partial_{n}, k_{n})$ et son dual}

Une seconde série numérotée, 1 à 3, reprend le calcul en toutes dimensions.
Pour $n \geq 1$, on se donne un $\mathfrak{S}_{n}$-ensemble $E_{n}$ et un
couple
\[
  \partial_{n} \colon E_{n} \longrightarrow E_{n+1}, \qquad
  k_{n} \colon E_{n+1} \longrightarrow E_{n},
\]
les noms venant des applications d'indices sous-jacentes : $\partial_{n}$
répond à l'inclusion $[1,n] \hookrightarrow [1,n+1]$ et $k_{n}$ à sa
rétraction. Les relations sont
\[
  k_{n}\partial_{n} = \mathrm{id}_{E_{n}}, \qquad
  \partial_{n}\sigma = i_{n}(\sigma)\,\partial_{n} \ \ (\sigma \in \mathfrak{S}_{n}),
  \qquad
  \sigma k_{n} = k_{n}\,i'_{n}(\sigma) \ \ (\sigma \in \mathfrak{S}_{n-1}),
  \qquad
  k_{n}\sigma_{n+1} = k_{n},
\]
où $i_{n} \colon \mathfrak{S}_{n} \hookrightarrow \mathfrak{S}_{n+1}$ est
l'inclusion, $i'_{n} = i_{n}i_{n-1}$, et $\sigma_{m}$ la transposition
$(m-1,m)$. La composée dans l'autre sens,
\[
  \rho_{n+1} = \partial_{n}k_{n} \colon E_{n+1} \longrightarrow E_{n+1},
\]
est alors un idempotent, projecteur de $E_{n+1}$ sur l'image de $E_{n}$ : c'est
la seule chose que $k_{n}\partial_{n} = \mathrm{id}$ permette d'affirmer, et
elle suffit.

\subsection*{Le système dual, et la condition de concaténation}

Renversant toutes les flèches, on obtient $E^{*}_{n}$ avec
$d_{n} = k^{*}_{n} \colon E^{*}_{n} \to E^{*}_{n+1}$ et
$p_{n} = \partial^{*}_{n} \colon E^{*}_{n+1} \to E^{*}_{n}$, satisfaisant les
relations transposées. Deux familles d'applications de descente y sont
définies, et il faut les distinguer soigneusement, car c'est leur \emph{couple}
qui porte tout l'énoncé :
\[
  \alpha^{*}_{n,p} = \partial^{*}_{p}\,\partial^{*}_{p+1}\cdots\partial^{*}_{n-1}
  \colon E^{*}_{n} \longrightarrow E^{*}_{p} \qquad (n > p \geq 1),
\]
\[
  \beta^{*}_{n,q} = \partial'^{*}_{q}\,\partial'^{*}_{q+1}\cdots\partial'^{*}_{n-1}
  \colon E^{*}_{n} \longrightarrow E^{*}_{q} \qquad (n \geq q \geq 1),
  \qquad \partial'_{m} = \rho_{m}\,\partial_{m},
\]
$\rho_{m}$ étant la permutation circulaire de $\mathfrak{S}_{m+1}$. Au niveau
des indices, $\alpha_{n,p}$ est l'inclusion du segment initial
$[1,p] \subset [1,n]$ et $\beta_{n,q}$ celle du segment final, $i \mapsto
i + (n-q)$ : $\alpha^{*}$ est donc la restriction au début et $\beta^{*}$ la
restriction à la fin.

La condition que le système doit satisfaire, en plus des relations, est alors
\[
  (**) \qquad
  E^{*}_{p+q} \xrightarrow{\ (\alpha^{*}_{p+q,p},\ \beta^{*}_{p+q,q})\ }
  E^{*}_{p} \times E^{*}_{q} \quad \text{surjectif}, \qquad p,q \geq 1 :
\]
un objet de longueur $p$ et un objet de longueur $q$ étant donnés, il existe
toujours un objet de longueur $p+q$ dont ils sont respectivement le début et la
fin. C'est une condition de concaténation, et c'est le dessin de la page 17 —
deux segments mis bout à bout, avec le point de raccord
marqué.\footnote{Elle a la forme des conditions de remplissage de la théorie
simpliciale, mais on notera que le but est un \emph{produit} et non un produit
fibré : les deux segments d'indices $[1,p]$ et $[p+1,p+q]$ sont disjoints, il
n'y a rien sur quoi les recoller. C'est la même forme que Hyp 2 de la section
précédente, où la famille des $U_{\gamma}$ devait couvrir un produit et non un
produit fibré.}

\subsection*{Le même système écrit en trois termes}

Les pages 18 à 20 redescendent au cas de trois étages et l'écrivent
$X_{1}, X_{2}, X_{3}$, avec $p_{0} \colon X_{2} \to X_{1}$,
$p_{12} \colon X_{3} \to X_{2}$, $\delta_{0} \colon X_{1} \to X_{2}$,
$\delta_{1},\delta_{2},\delta_{3} \colon X_{2} \to X_{3}$ et les groupes
$\mathfrak{S}_{2}$, $\mathfrak{S}_{3}$. C'est le système
$E_{0}, E_{1}, E_{2}$ des pages 2 à 8, décalé d'un
cran.\footnote{Ces $X_{n}$ n'ont rien à voir avec les $X_{i}$ de
« Structures stratifiées », qui sont les parties fermées d'une stratification.
La collision de lettres est réelle et court sur tout le dossier ; nous gardons
ici $E_{0}, E_{1}, E_{2}$ pour le système à trois étages, et réservons $X$ à
l'espace stratifié.}

Deux carrés y sont déclarés cartésiens par l'auteur — il écrit « cart » en leur
centre —, et ils sont ensuite récrits avec $E$ au lieu de $X$ et $\partial$ au
lieu de $\delta$ : le passage d'un système à l'autre est ce que la page cherche.
Le second de ces carrés,
\begin{tikzcd}[column sep=large, row sep=small]
E_{1} & E_{2} \arrow[l, "p_{12}"'] \\
E_{0} \arrow[u, "\delta_{0}"] & E_{1} \arrow[l, "p_{0}"] \arrow[u, "\delta_{1}"']
\end{tikzcd}

dit exactement, une fois $E_{0}$ réduit à un point, que l'image de $E_{1}$ dans
$E_{2}$ est la fibre de $p_{12}$ au-dessus du point marqué. C'est cette
assertion que le lot suivant démontre, à la page 24.

La page 19 pose alors l'hypothèse qui ouvre le lot suivant : $X_{1}$ réduit à un
élément, donc $E_{1}$ un $\mathfrak{S}_{2}$-ensemble pointé et $E_{2}$ un
$\mathfrak{S}_{3}$-ensemble pointé, reliés par
\[
  d \colon E_{1} \rightleftarrows E_{2} \colon p, \qquad p\,d = \mathrm{id}_{E_{1}}.
\]
La page 20 ne porte que trois lignes de calcul en coordonnées, qui reprennent
$d$ et $\rho^{2}d$ ; elles sont reprises et menées à leur terme à la page
21.\footnote{Les deux triplets de la page 20 sont donnés par la transcription
comme des lectures incertaines. Le calcul de la page 21, où les mêmes formules
sont écrites lisiblement, donne $d(x,y) = (x,y,y)$ et
$\rho^{2}d(x,y) = (y,x,y)$ ; c'est cette lecture-là que nous retenons.}

\subsection*{Une page intercalée}

La page 17, écrite au dos de la même lettre administrative que ses voisines,
n'appartient pas à cette série mais au texte « Hyperrecouvrements » : elle
porte le calcul de $\varphi_{!}(F \times G)$ par colimites qui justifie le
passage des représentables aux objets quelconques, et le couple adjoint
$(\varphi_{!}, \varphi^{*})$ écrit en toutes lettres — « $u$ continu, $v$
cocontinu ». Nous l'avons rendue à sa place logique, ci-dessus.

\end{document}
